%% ============================================================
%% MODELO DE ATIVIDADE · FIRMWARE E SOFTWARE
%%
%% Copie este arquivo para atividades/ (atividade-02.tex, por exemplo),
%% inclua-o no relatorio-aacc.tex com \input{atividades/atividade-02}
%% e escreva. Cada \preencher{…} é uma pendência: troque-o pelo seu
%% texto. Cada \figuraprovisoria também: troque-a pela figura de verdade.
%%
%% Uma atividade é UMA entrega sua, com começo, meio e fim. "Participei
%% do projeto X por dois anos" não é uma atividade: são várias. Divida.
%% ============================================================
\begin{atividade}{Módulo M do firmware: requisitos, arquitetura e testes}{
  tipo         = software,
  modalidade   = {},   % a chave da modalidade pedida: EXT, PES, COMP, VPC, EVT, PRD, ORG, ENS, GES, CUR, OUT
  alternativa  = {},   % a modalidade alternativa, se o Colegiado não aceitar a primeira
  horas        = {},   % as horas desta atividade, com base nos registros (Parte IV)
  inicio       = {},   % AAAA-MM
  fim          = {},   % AAAA-MM, ou: em andamento
  projeto      = {},
  papel        = {},   % Responsável, Corresponsável ou Colaborador(a)
  supervisor   = {},   % quem confirma: nome e cargo (o supervisor do projeto, o gerente)
  registros    = {},   % os códigos no SOMA, com ponto e vírgula: NBL-14; NRO-PRO-003-7
  competencias = {},   % as competências das DCN que ela desenvolveu (I a VIII): I, III, V
  disciplinas  = {},   % as disciplinas do seu curso ligadas a ela: ELE075 Eletrônica; ELE067 Circuitos
}

\begin{contexto}
\preencher{Que sistema, que problema, e com que requisitos (tempo real, memória, consumo, segurança).}
\end{contexto}

\begin{contribuicao}
\preencher{Os módulos que você escreveu, os PRs que abriu e revisou (com os números), as decisões que tomou.}
\end{contribuicao}

\begin{desenvolvimento}
\fichatecnica{
  linguagens    = {},   % Linguagens
  plataforma    = {},   % Plataforma (MCU, sistema, framework)
  repositorio   = {},   % Repositório
  contribuicoes = {},   % Contribuições (commits, PRs, revisões)
  testes        = {},   % Testes
  licenca       = {},   % Licença — opcional
}

\preencher{A arquitetura: os módulos, as interfaces e o fluxo de dados. As decisões registradas (ADR, NRO-PRO-013). Os trechos que importam, explicados --- não o código inteiro. Os testes.}

\begin{figure}[htbp]
  \centering
  \figuraprovisoria{A arquitetura: os módulos e o fluxo de dados (ou uma listagem com legenda, que também conta).}
  % \includegraphics[width=0.9\linewidth]{figuras/arquitetura.pdf}
  \caption{Arquitetura do firmware do módulo M.}
  \fonte{Elaborado pelo autor.}
  \label{fig:arquitetura}
\end{figure}

% Um trecho de código que importa, com legenda (conta como figura):
% \begin{lstlisting}[language=C, caption={Rotina de interrupção do conversor A/D.}, label={lst:isr}]
% void ADC_IRQHandler(void) { ... }
% \end{lstlisting}
\end{desenvolvimento}

\begin{resultados}
\preencher{As métricas: latência, taxa de amostragem, consumo, memória, falhas. Os testes que passam e o que foi para a versão em uso.}
\end{resultados}

\begin{evidencias}
% Uma linha por evidência: \evidencia{tipo}{identificador}{o que mostra}{onde conferir}
% \evidencia{Repositório}{github.com/neurodynamics-dev/<repo>, PR \#<n>}{O pull request do módulo, com a revisão}{GitHub (acesso a pedido)}
% \evidencia{Arquivo controlado}{NRO-PRO-013-<PN>}{O registro da decisão de arquitetura}{SOMA › Arquivos}
\end{evidencias}

\begin{aprendizados}
\preencher{O que a programação de um sistema real ensinou que as disciplinas de programação e de sistemas embarcados não ensinaram.}
\end{aprendizados}
\end{atividade}
